\section*{第 \ref{ch:g10:cell-metabolism} 章——细胞代谢}

\begin{solution}{exo:g10:cell-metabolism:1}
\omterm{def:g10:cell-metabolism:metabolism}{代谢}是\omterm{def:g10:cells-common-unit:cell}{细胞}内所有化学反应的总和，包括分解分子以获取能量，
以及合成\omterm{def:g10:cells-common-unit:cell}{细胞}自身的物质。酶是一种\omterm{def:g10:chemistry-of-life:families}{蛋白质}，能加快
某一特定反应，它与底物结合，
并释放出产物，自身不被消耗。
\end{solution}

\begin{solution}{exo:g10:cell-metabolism:2}
呼吸作用：$\mathrm{C_6H_{12}O_6} + 6\,\mathrm{O_2} \to
6\,\mathrm{CO_2} + 6\,\mathrm{H_2O}$ + 能量。\omterm{def:g10:cell-metabolism:autotrophy}{光合作用}：
$6\,\mathrm{CO_2} + 6\,\mathrm{H_2O}$ + 光 $\to \mathrm{C_6H_{12}O_6}
+ 6\,\mathrm{O_2}$。物质的总体平衡方向相反，但
两者的反应、酶、场所（\omterm{def:g10:cells-common-unit:organelle}{线粒体}与\omterm{def:g10:cells-common-unit:organelle}{叶绿体}）以及
能量来源（化学键与光）完全不同。
\end{solution}

\begin{solution}{exo:g10:cell-metabolism:3}
\omterm{def:g10:cell-metabolism:autotrophy}{异养}生物：蘑菇、胡萝卜根\omterm{def:g10:cells-common-unit:cell}{细胞}、酸奶中的细菌、人的
肌肉\omterm{def:g10:cells-common-unit:cell}{细胞}。\omterm{def:g10:cell-metabolism:autotrophy}{自养}生物：苔藓、小球藻。
\end{solution}

\begin{solution}{exo:g10:cell-metabolism:4}
二氧化碳使气球膨胀；乙醇在
液体中积累。烧瓶中缺氧：若通入空气，酵母
会进行呼吸作用，不产生乙醇，且消耗的糖少得多。
\end{solution}

\begin{solution}{exo:g10:cell-metabolism:5}
过氧化氢酶是一种\omterm{def:g10:chemistry-of-life:families}{蛋白质}；沸腾会破坏其结构，从而
使其失去活性，而过氧化氢本身不会迅速分解。
\end{solution}

\begin{solution}{exo:g10:cell-metabolism:6}
$(7.9 - 0.9)/7 = \qty{1.0}{mg/L}$ 每分钟。曲线趋于平缓，
是因为氧气已耗尽；酵母随后转为\omterm{def:g10:cell-metabolism:fermentation}{发酵}，
而\omterm{def:g10:cell-metabolism:fermentation}{发酵}不消耗氧气。
\end{solution}

\begin{solution}{exo:g10:cell-metabolism:7}
$0.8 + 0.2 = \qty{1.0}{mg/L}$ 每分钟：光照下的斜率为
产生量减去光照下仍在进行的呼吸作用。
\end{solution}

\begin{solution}{exo:g10:cell-metabolism:8}
\qty{18}{g} 葡萄糖为 \qty{0.1}{mol}，需要 \qty{0.6}{mol}
氧气，$0.6 \times 32 = \qty{19.2}{g}$，并释放 \qty{0.6}{mol}
二氧化碳，$0.6 \times 44 = \qty{26.4}{g}$。
\end{solution}

\begin{solution}{exo:g10:cell-metabolism:9}
萌发的豌豆分解储备物质进行呼吸作用，并释放热量；
煮过的豌豆已经死亡，其中的酶被破坏，因此不发生反应，也不释放热量。若密封不透气，
氧气会耗尽，呼吸作用随之停止。
\end{solution}

\begin{solution}{exo:g10:cell-metabolism:10}
暴露在光下的一半会被染成蓝黑色（淀粉由\omterm{def:g10:cell-metabolism:autotrophy}{光合作用}产生）；
被遮盖的一半则不会。48 小时的黑暗使叶片
耗尽原有的淀粉，因而检测到的淀粉
是当天光照的结果。
\end{solution}

\begin{solution}{exo:g10:cell-metabolism:11}
\omterm{def:g10:cell-metabolism:fermentation}{发酵}每分解一分子葡萄糖所产生的 ATP，约为呼吸作用的
十五分之一；为获得所需的 ATP，\omterm{def:g10:cell-metabolism:fermentation}{发酵}\omterm{def:g10:cells-common-unit:cell}{细胞}每分钟
必须多分解十五倍的糖。
\end{solution}

\begin{solution}{exo:g10:cell-metabolism:12}
比赛期间，肌肉对 ATP 的需求超过了血液所能输送的氧气；
\omterm{def:g10:cells-common-unit:cell}{细胞}转而进行乳酸\omterm{def:g10:cell-metabolism:fermentation}{发酵}，其产物乳酸随之积累。
慢跑时，氧气供应能够跟上，仅靠呼吸作用即可提供 ATP，
不会产生乳酸。
\end{solution}

\begin{solution}{exo:g10:cell-metabolism:13}
环境发生了变化（无光照，有机食物）；\omterm{def:g10:universal-dna:gene}{基因}并未改变。
后代重新形成\omterm{def:g10:cells-common-unit:organelle}{叶绿体}，因此相关的\omterm{def:g10:universal-dna:information}{遗传信息}从未丢失：
丢失的只是正在使用的结构，而非
\omterm{def:g10:universal-dna:gene}{基因}所允许的各种可能性。
\end{solution}

\begin{solution}{exo:g10:cell-metabolism:14}
植物在光下进行\omterm{def:g10:cell-metabolism:autotrophy}{光合作用}，提供氧气和\omterm{def:g10:chemistry-of-life:mineral-organic}{有机物}（即蜗牛所取食的叶子）；
蜗牛进行呼吸作用，为植物提供二氧化碳。在黑暗中，植物只进行呼吸作用：
二者都消耗氧气，却无人产生氧气，
几天之内两者都会死亡。
\end{solution}

\begin{solution}{exo:g10:cell-metabolism:15}
每分钟分解的葡萄糖为原来的 $30/2 = 15$ 倍。酿酒者需要乙醇，
而只有\omterm{def:g10:cell-metabolism:fermentation}{发酵}才能产生乙醇；进行呼吸作用的酵母
会把糖转化为二氧化碳、水和更多的酵母，而不产生酒精。
\end{solution}

\begin{solution}{pb:g10:cell-metabolism:1}
\textbf{1.} 酒精\omterm{def:g10:cell-metabolism:fermentation}{发酵}，因为密闭的\omterm{def:g10:cell-metabolism:fermentation}{发酵}罐中不含
氧气：$\mathrm{C_6H_{12}O_6} \to 2\,\mathrm{C_2H_5OH} +
2\,\mathrm{CO_2}$。

\textbf{2.} $200/180 \approx \qty{1.11}{mol}$。

\textbf{3.} 乙醇 $2 \times 1.11 = \qty{2.22}{mol}$，即 $2.22
\times 46 \approx \qty{102}{g}$；体积 $102/0.79 \approx
\qty{129}{mL}$；约为 13\%（按体积计）。

\textbf{4.} $2.22 \times 44 \approx \qty{98}{g}$ 的二氧化碳；
$2.22 \times 24 \approx \qty{53}{L}$ 的气体。

\textbf{5.} 约 \qty{53000}{L} $= \qty{53}{m^3}$。酒窖的容积为
$4 \times 6 \times 2.5 = \qty{60}{m^3}$：释放出的气体
几乎可充满它。

\textbf{6.} $\mathrm{C_6H_{12}O_6} + 6\,\mathrm{O_2} \to
6\,\mathrm{CO_2} + 6\,\mathrm{H_2O}$。

\textbf{7.} 氧气 $6 \times 1.11 = \qty{6.67}{mol}$，$6.67 \times 32
\approx \qty{213}{g}$；二氧化碳 $6.67 \times 44 \approx
\qty{293}{g}$。

\textbf{8.} 三倍：呼吸作用将葡萄糖中的全部六个碳
原子都转化为二氧化碳，\omterm{def:g10:cell-metabolism:fermentation}{发酵}只转化两个，其余四个
留在乙醇中。

\textbf{9.} 有氧呼吸时，持续时间长十五倍：每分子葡萄糖
可多产生十五倍的 ATP。

\textbf{10.} 乙醇只能通过\omterm{def:g10:cell-metabolism:fermentation}{发酵}产生；若是通气的\omterm{def:g10:cell-metabolism:fermentation}{发酵}罐，
则只会得到二氧化碳、水和酵母，而没有葡萄酒。

\textbf{11.} $1.11 \times 70 \approx \qty{78}{kJ}$，每升；
即整个\omterm{def:g10:cell-metabolism:fermentation}{发酵}罐 \qty{78000}{kJ}，亦即 \qty{78}{MJ}。

\textbf{12.} $78/4.2 \approx \qty{18.5}{\celsius}$。

\textbf{13.} $20 + 18.5 \approx \qty{38.5}{\celsius}$：接近酵母的耐受极限，
温暖的酒窖会将其推过极限。大型\omterm{def:g10:cell-metabolism:fermentation}{发酵}罐
因表面积相对于体积较小而散热缓慢，
因此必须人为降温。

\textbf{14.} 酶是\omterm{def:g10:chemistry-of-life:families}{蛋白质}；高温使其变性，
酶失去作用。没有酶，酵母便无法转化糖，
无论剩余多少糖。

\textbf{15.} 蕴藏于乙醇的化学键中：每摩尔葡萄糖约
\qty{2800}{kJ}，这正是乙醇能够燃烧的原因。

\textbf{16.} 在地面上，形成一层，并随\omterm{def:g10:cell-metabolism:fermentation}{发酵}进行而增厚；
低处拿着的蜡烛会在二氧化碳已取代氧气的地方熄灭，
——这是人受影响之前的警示。

\textbf{17.} 乳酸\omterm{def:g10:cell-metabolism:fermentation}{发酵}。其产生的 ATP 少十五倍；
大脑和心脏无法依靠这种方式满足其需要，乳酸便
积累；一两分钟内便会出现意识丧失。

\textbf{18.} 高浓度的乙醇会损伤酵母的
膜和酶：\omterm{def:g10:cells-common-unit:cell}{细胞}自身的产物毒害了
它们的\omterm{def:g10:cell-metabolism:metabolism}{代谢}。

\textbf{19.} 气泡是\omterm{def:g10:cell-metabolism:fermentation}{发酵}产生的二氧化碳。乙醇
在烘箱中蒸发（沸点为 \qty{78}{\celsius}）。

\textbf{20.} 按体积计约 13\% 的酒精，以及每升葡萄汁
产生 \qty{53}{L} 的气体；两者都源于酵母在无氧条件下，
将每个葡萄糖分子\omterm{def:g10:cell-metabolism:fermentation}{发酵}为两个乙醇和两个二氧化碳。
\end{solution}
